% ======================================================================
% SECTION 4: HR 4502 - Cancer Patient Robot-and-Humanoid Choice Act of 2026
% Revision of California AB 2847 (2026) and FDA AI Regulatory
% Decision-Making Draft Guidance 2025.
% ======================================================================

[PROCESSING INSTRUCTION (HR 4502 - Cancer Patient Robot-and-Humanoid
Choice Act of 2026, 6 to 9 paragraphs structured into the five
legislative-act subsections):

The bill must explicitly state in its first paragraph that it is a
"Revision of California AB 2847 (2026) Patient Rights and Robotic
Safety Act and FDA AI Regulatory Decision-Making Draft Guidance 2025"
so that readers can trace the chain of prior authority. The CA AB 2847
prior text is documented at
new-trial/site/02-legislation-patient-rights/physical_ai_patient_rights.tex
and at national-platform/new_trial_psl/02_ab2847_patient_rights.tex
(read both files in full when populating). The FDA AI Regulatory
Decision-Making Draft is at \cite{fda_ai_regulatory_decision_draft_2025}.

Subsection 4.1 - Findings and Declarations: Read
patients/patient_robot_instructions_fixed.tex pages 1 through 10
(also chunked at
patients/patient_robot_instructions_fixed_chunk/ via
01_preamble_robots1to5.tex and 02_robots6to10_closing.tex) and
quote the 10 robot categories enumerated there: surgical robots,
collaborative cobots, RT positioning robots, robotic needle-placement
systems, social companion robots, humanoids, RT motion-management
robots, imaging assistant robots, steerable needle robots,
rehabilitation exoskeletons. Read
national-platform/usl_standard/01_preamble_introduction_framework.tex
and 02_results_discussion_closing.tex for the USL scoring framework.
Compose 5 to 7 enumerated findings citing
\cite{kawchak_2026_18810541}, \cite{kawchak_2026_18778219},
\cite{target2025robotic_bronchoscopy}, and
\cite{vorontsov2024virchow}. At least one finding must reference the
California AB 2847 prior bill being revised. At least one finding must
reference the FDA AI Draft 2025 capability-neutrality requirement.

Subsection 4.2 - Definitions: Define the operational terms used
throughout the bill:
- "Patient robot choice" means the patient's right to select which
  surgical robot, humanoid robot, or companion robot performs their
  care. Reference patients/patient_robot_instructions_fixed.tex
  pages 1 to 10.
- "Qualified robot" means a Physical AI system that meets the
  Unification Standard Level (USL) threshold for its category, as
  documented in national-platform/usl_standard/ and
  unification/usl/paper/usl_oncology_trials.tex. Cite
  \cite{kawchak_2026_18778219}.
- "Robot category" means one of the 10 categories enumerated in the
  prior California SB 1042 Section 2 definitions, repeated at
  new-trial/site/01-legislation-authorization/physical_ai_trial_authorization.tex
  and at national-platform/new_trial_psl/01_preamble_and_sb1042.tex,
  with each category carrying its own USL minimum threshold. Cite
  \cite{kawchak_2026_19176370}.
- "Capability-neutral" means the bill's definitions must apply to any
  robot architecture (LLM-driven, foundation-model, agentic system,
  or specialized hardware) that meets the USL threshold, drawing on
  patients/paper/Deep-Research-3/part2_ai_robotics_legislation.md
  guardrail #2 and \cite{fda_ai_regulatory_decision_draft_2025}.

Subsection 4.3 - Patient Robot Choice Rights (operative core):
Compose 5 to 7 enumerated rights:
- Right to receive a side-by-side comparison of all qualified robots
  for the patient's procedure, including USL scores, prior outcomes,
  and named instances. Use the named-instance examples from
  patient-journey/stage_05_surgery.py (da Vinci Xi at USL 87.5 for
  lobectomy), stage_06_recovery.py (post-op cobot for limb stabili-
  zation), stage_07_immunotherapy.py (Franka Emika at USL 88.75 for
  pembrolizumab), and stage_09_surveillance.py (companion robot for
  pediatric or geriatric patients). Cite \cite{kawchak_2026_19119939}.
- Right to select the specific robot instance (e.g., RTPOS-01 vs
  RTPOS-03) following the patient_records.md naming pattern at
  new-trial/national-24-7-trial/hour-12/patient_records.md.
- Right to substitute a humanoid robot (e.g., Tesla Optimus, Boston
  Dynamics Atlas, Agility Digit, evaluated under
  unification/usl/humanoids/) for routine non-surgical tasks
  (transport, vitals collection, medication delivery) where the
  patient prefers reduced human contact.
- Right to refuse a specific robot instance and request a substitute
  at no additional cost to the patient.
- Right to receive a USL transparency report (the four-dimension USL
  scoring sheet) for any robot before consenting to its use. Cite
  \cite{kawchak_2026_18778219}.
- Right of appeal to the FDA Oncology Center of Excellence DCT
  pathway for any robot denial that the patient believes is
  capability-equivalent. Cite \cite{fda_oce_advancing_oncology_dct_2024}.

Subsection 4.4 - Implementation: 12-month deadline for ALL qualified
sites to publish a per-robot USL report; FDA, AHRQ, and ONC roles;
explicit reference that this is a Revision of CA AB 2847 (2026) for
NATIONAL coverage; preemption rule preserves stronger state
protections. Read national-platform/patient_robot/02_robots6to10_closing.tex
for the operational implementation pattern and
sponsor/final_paper/scripts/safety/ for the safety-workflow integration.

Subsection 4.5 - Reporting and Enforcement: Annual federal report
listing the per-state, per-cancer, per-robot-category selection rate;
the AI-decision-tracking infrastructure at
sponsor/final_paper/scripts/core_agents/ (53 agents, see
sponsor/final_paper/scripts/core_agents/__init__.py for the list) that
the FDA can audit; civil penalties harmonized with the FDA AI Draft
2025 capability-neutrality enforcement framework.

Close the bill (1 paragraph after Subsection 4.5) with two patient-
control framings:
(i) FOR INDIVIDUAL PATIENTS: A 58-year-old female NSCLC patient
    matching the PAT-2026-0042 profile in patient-journey/ chooses the
    da Vinci Xi for her lobectomy after reviewing USL 87.5 vs the
    USL 84.0 alternative, and chooses a Tesla Optimus humanoid for
    home-based vitals collection during her 35 pembrolizumab cycles
    at \cite{dronca2026cancer_care_beyond_walls}.
(ii) FOR BROAD ADOPTION: All Physical AI oncology trial sites must
     publish a USL transparency report for every deployed robot
     instance within 12 months. The bill positions the United States
     as Number 1 in the world for patient-controlled robot selection,
     replacing the prior provider-only choice pattern under
     traditional CA AB 2847 with a national patient-led standard.
     Cite \cite{kawchak_2026_18810541} and \cite{kawchak_2026_19176370}.]
